\documentclass[12pt]{article}

\usepackage{setspace}
\onehalfspace

%\usepackage[utf8]{inputenc}
%\usepackage{t1enc}
%\usepackage[T1]{fontenc}

\usepackage{amsmath}

\usepackage{moodle}

\DeclareMathOperator{\Exp}{Exp}
\DeclareMathOperator{\Norm}{{\cal N}}
\DeclareMathOperator{\Binom}{Binom}
\DeclareMathOperator{\E}{E}
\DeclareMathOperator{\D}{D}
\renewcommand{\P}{\operatorname{P}}
%\DeclareMathOperator{\Pr}{P}
\DeclareMathOperator{\F}{F}
\DeclareMathOperator{\R}{\mathbb{R}}
%\DeclareMathOperator{\Norm}{{\cal N}}
\DeclareMathOperator{\Unif}{{\cal U}}


\DeclareMathOperator{\cov}{cov}





\begin{document}

\begin{quiz}{gyakorló}

\begin{multi}{val7}
Bármely $C$-re $\P(C)=\sum_{\omega_k\in C} \P( \{\omega_k \})$.
\item* igaz
\item hamis
\end{multi}
\begin{multi}{dvv8}
Egy $\xi$ valségi változó az $x_1,\ldots x_n$ értékeket veheti fel és $p_k=\P(\xi=x_k)$. 
Ekkor a második momentuma:
$$
\E(\xi^2)=\sum_{k=1}^n x_k^2 p_k
$$
\item* igaz
\item hamis
\end{multi}
\begin{multi}{felt2}
$A$-nak a $B$-re vonatkozó feltételes valsége $\P(A|B)=\frac{\P(A)}{\P(B)}$, ha $A$ maga után vonja $B$-t.
\item* igaz
\item hamis
\end{multi}
\begin{multi}{impfvv3}
Egy $\xi\sim \Unif(a,b)$ szórásnégyzete: $\frac{(b-a)^2}{12}$
\item* igaz
\item hamis
\end{multi}
\begin{multi}{impdvv2}
Egy $\xi\sim \Binom(n,p)$ várható értéke: $n+p$
\item igaz
\item* hamis
\end{multi}
\begin{multi}{felt8}
Ha $A_1,\ldots,A_n$ teljes eseményrendszer és $\P(B)>0$ akkor:
$$
\P(B)=\sum_{k=1}^n \P(A_k|B)\P(B)
$$
\item* igaz
\item hamis
\end{multi}
\begin{multi}{impfvv7}
Ha $\xi\sim \Norm(\mu,\sigma^2)$ akkor $\frac{\xi-\mu}{\sigma}\sim \Norm(0,1)$.
\item* igaz
\item hamis
\end{multi}
\begin{multi}{esem7}
Az $A$ és $B$ események pontosan akkor függetlenek, ha $\P(AB)=\P(A)\P(B)$.
\item* igaz
\item hamis
\end{multi}
\begin{multi}{eofv2}
Egy $\F:\R\to [0,1]$ pontosan akkor eloszlásfüggvény, ha monoton nemcsökkenő, 
$\lim_{x\to+\infty}\F(x)=1$ és $\lim_{x\to-\infty}\F(x)=0$
\item igaz
\item* hamis
\end{multi}
\begin{multi}{val1}
A valószínűség additív, azaz $\P(A+B)=\P(A)+\P(B)$.
\item igaz
\item* hamis
\end{multi}
\begin{multi}{dvv16}
$\xi$ és $\eta$ véges értékkészletű valségi változókra:
$$
\D^2(\xi+\eta)=\D^2(\xi)+\D^2(\eta)
$$
\item igaz
\item* hamis
\end{multi}
\begin{multi}{eofv10}
Bármely $\xi$-re és $a$-ra $\P(\xi= a)=0$, feltéve hogy $\F_{\xi}$ folytonos.
\item* igaz
\item hamis
\end{multi}
\begin{multi}{esem8}
Az $A$ és $B$ események pont akkor függetlenek, ha $\P(A+B)=\P(A)+\P(B)$.
\item igaz
\item* hamis
\end{multi}
\begin{multi}{dvv17}
$\xi$ és $\eta$ $\xi$ és $\eta$ véges értékkészletű valségi változókra:
$$
\D^2(\xi+\eta)=\D^2(\xi)+\D^2(\eta)+2\cov(\xi,\eta)
$$
\item* igaz
\item hamis
\end{multi}
\begin{multi}{esem5}
Ha $\P(A)=1$ akkor $A$ bekövetkezhet.
\item* igaz
\item hamis
\end{multi}
\begin{multi}{impdvv4}
Egy $\xi\sim \Binom(n,p)$ szórásnégyzete: $\frac{n+1}{2}$
\item igaz
\item* hamis
\end{multi}
\begin{multi}{felt10}
A teljes függetlenségből következik a páronkénti.
\item* igaz
\item hamis
\end{multi}
\begin{multi}{dvv14}
$\xi$ és $\eta$ független, véges értékkészletű valségi változókra:
$$
\D^2(\xi+\eta)=\D^2(\xi)+\D^2(\eta)
$$
\item* igaz
\item hamis
\end{multi}
\begin{multi}{eofv8}
Bármely $\xi$-re és $a$-ra $\F_{\xi}(a)+\P(\xi\ge a)=1$.
\item* igaz
\item hamis
\end{multi}
\end{quiz}

\end{document}

